\section{Small spectrum and rescaled linearized evolution}\label{sec:spectrum}

This section concerns the HYM linearization at the constant flat metric,
and therefore the scalar Laplacian on matrix entries. It supplies an
operator identification and uniform Hilbert-space estimates. It makes no
assertion that a nonstationary metric has the same instantaneous spectrum.

Let $V=\RR^N$ with inner product
$\langle a,b\rangle_A=\sum_iA_i a_i b_i$. Define
\[
 q(a)=\sum_{p:\,i\leftrightarrow j}(a_i-a_j)^2,
 \qquad B=2\pi A^{-1}Q,
\]
where $Q$ is the symmetric matrix of $q$. The component graph is connected,
so its kernel consists of constant vectors. Write
$0=\mu_0<\mu_1\le\cdots\le\mu_{N-1}$ for the eigenvalues of
$A^{-1}Q$. Let $0=\lambda_0(t)<\lambda_1(t)\le\cdots$ be the
scalar fibre spectrum, counted with multiplicity, and let $S_t$ be the
span of the first $N$ eigenfunctions.

Let $\Omega_i(r,t)$ denote the image under Lemma~\ref{lem:plumbing}
of $C_i$ with the discs of coordinate radius $r$ about all its node
punctures removed. Here $r$ is fixed, positive, and smaller than the
chosen outer plumbing radius. In this section $o_r(1)$ denotes a quantity
tending to zero as $t\to0$ with $r$ fixed.

We use the following scalar facts from \cite{wuebben-curves}. They hold
for the present induced K\"ahler metrics, and indeed under weaker
hypotheses in that paper:
\begin{equation}\label{eq:scalar-law}
 L\lambda_j(t)\longrightarrow2\pi\mu_j\quad(0\le j<N),
 \qquad\lambda_N(t)\ge\gamma>0.
\end{equation}
We also need the compactness and form estimates in its proof, rather
than eigenvalues alone. For clarity we state precisely the consequences
used here. On every fixed thick component subset, a unit vector in $S_t$
differs in $L^2$ from its component average by $O_r(L^{-1/2})$.
Its squared mass in fixed node neighbourhoods of radius $r$ is at
most $Cr^2+o_r(1)$. These estimates hold uniformly for all combinations
of the $N$ eigenfunctions. Finally the capacitary lower bound gives
\begin{equation}\label{eq:scalar-form-lower}
 L\int_{X_t}|du|^2\,dA_t
 \ge2\pi q(a(u))-o_r(1)\|u\|^2\qquad(u\in S_t),
\end{equation}
where $a(u)$ consists of averages over $\Omega_i(r,t)$, using the
same fixed coordinate radius as the capacitary cutoffs.
In \cite[Section~4]{wuebben-curves} these statements
are the thick estimates, the node-mass estimate, their extension to
finite combinations, and the energy lower bound. The lower bound is
obtained by pairing $du$ with the differentials of logarithmic
capacitary potentials. Thus it applies to an entire low subspace,
not only to individual eigenfunctions.

Fix a sufficiently small outer radius $r_0$, and define
\[
 T_t:S_t\to V,\qquad (T_tu)_i=
 \frac{1}{\operatorname{area}\Omega_i(r_0,t)}
              \int_{\Omega_i(r_0,t)}u\,dA_t.
\]
Adjoints in what follows use the actual fibre $L^2$ inner product and
the weighted inner product on $V$.

\begin{theorem}\label{thm:operator}
For sufficiently small $t$, the averaging map $T_t$ is invertible.
The polar identification
\begin{equation}\label{eq:polar}
 J_t=T_t^*(T_tT_t^*)^{-1/2}:V\longrightarrow S_t
\end{equation}
is an isometry, and
\begin{equation}\label{eq:operator-limit}
 \|J_t^*L\Delta_tJ_t-B\|_{V\to V}\longrightarrow0.
\end{equation}
For every $T<\infty$,
\begin{equation}\label{eq:low-heat}
 \sup_{0\le\tau\le T}
 \|J_t^*e^{-L\tau\Delta_t}J_t-e^{-\tau B}\|\longrightarrow0.
\end{equation}
For $0<\delta<T$, regarding $J_t$ as an embedding into the full
fibre $L^2$ space,
\begin{equation}\label{eq:full-heat}
 \sup_{\delta\le\tau\le T}
 \|e^{-L\tau\Delta_t}-J_te^{-\tau B}J_t^*\|_{L^2\to L^2}
 \longrightarrow0.
\end{equation}
The shifted finite resolvents converge uniformly when their positive
shift ranges over a compact interval. On $S_t^\perp$,
\begin{equation}\label{eq:complement-bounds}
 \|\Delta_t^{-1}\|_{L^2\to L^2}\le\gamma^{-1},\qquad
 \|e^{-s\Delta_t}\|_{L^2\to L^2}\le e^{-\gamma s}.
\end{equation}
\end{theorem}
\begin{proof}
First we prove that $T_t$ is asymptotically isometric. Choose a fixed
smaller radius $r<r_0$. The thick Poincar\'e estimate implies that
the means at $r$ and $r_0$ differ by
$O_r(L^{-1/2})\|u\|$. On the thick pieces, the mass is the area
times the squared mean, up to $o_r(1)\|u\|^2$. The omitted mass is
bounded by $Cr^2\|u\|^2+o_r(1)\|u\|^2$, and the omitted component
area is $O(r^2)$. Uniformly on the unit sphere of $S_t$,
\[
 \bigl|\|T_tu\|_A^2-\|u\|^2\bigr|
 \le(Cr^2+o_r(1))\|u\|^2.
\]
Let $t\to0$ for fixed $r$ and then $r\to0$. This proves
$T_t^*T_t=I+o(1)$. Both spaces have dimension $N$, so $T_t$ is
invertible for small $t$; \eqref{eq:polar} is an isometry and
$T_tJ_t=(T_tT_t^*)^{1/2}\to I$.

Put $B_t=J_t^*L\Delta_tJ_t$. Use a fixed capacity radius compatible
with $r_0$ in \eqref{eq:scalar-form-lower}; changing between two
fixed averaging radii only adds $O(L^{-1/2})\|u\|$ to the means.
Substituting $u=J_ta$ gives, uniformly on $V$,
\[
 \langle B_ta,a\rangle_A
 \ge\langle Ba,a\rangle_A-\varepsilon_t\|a\|_A^2,
 \qquad\varepsilon_t\to0.
\]
The eigenvalues of $B_t$ are $L\lambda_j(t)$ for $0\le j<N$.
Equation~\eqref{eq:scalar-law} therefore gives
$\tr(B_t-B)\to0$. For the self-adjoint difference $D_t=B_t-B$,
\[
 \lambda_{\min}(D_t)\ge-\varepsilon_t,\qquad
 \lambda_{\max}(D_t)\le\tr D_t+(N-1)\varepsilon_t.
\]
This proves \eqref{eq:operator-limit}, including when limiting
eigenvalues are repeated. The form inequality is the information
which makes this inference stronger than convergence of spectra.

Since $S_t$ is a spectral subspace,
$J_t^*e^{-L\tau\Delta_t}J_t=e^{-\tau B_t}$.
Duhamel's formula for nonnegative self-adjoint matrices gives
\[
 \|e^{-\tau B_t}-e^{-\tau B}\|\le\tau\|B_t-B\|.
\]
The resolvent identity similarly gives the bound
$z^{-2}\|B_t-B\|$ for shifts $z>0$.
The spectral gap in \eqref{eq:scalar-law} proves
\eqref{eq:complement-bounds}. Decomposing the full fibre space into
$S_t\oplus S_t^\perp$ bounds the difference in \eqref{eq:full-heat}
by $T\|B_t-B\|+e^{-\gamma L\delta}$.
\end{proof}

The exclusion of $\tau=0$ in \eqref{eq:full-heat} is necessary:
the identity on the orthogonal complement cannot disappear in
operator norm at time zero. The polar identification need not carry
the finite constant vector to an exactly constant function for
each nonzero $t$; it does so asymptotically. No unshifted inverse
on a space containing a kernel is asserted.

Let $\mathfrak p_r$ be the real space of trace-free Hermitian matrices,
of dimension $r^2-1$. Tensoring Theorem~\ref{thm:operator} with its
identity proves all corresponding HYM linearization statements.
The finite equation linearized at $H_i=I$ is
\[
 A_i\partial_\tau u_i=2\pi\sum_{p:\,i\leftrightarrow j}(u_j-u_i),
\]
which is exactly $-B\otimes I_{\mathfrak p_r}$. The estimates
\eqref{eq:complement-bounds} are Hilbert-space inverse and heat
bounds, not global unweighted higher-regularity estimates. The
nonlinear proof instead uses Lemma~\ref{lem:comparison} and the
fixed-chart scalar mean-value inequality.
